\documentclass[10pt,a4paper]{amsart}
\usepackage{amsmath,amssymb,amsthm,graphicx}
% Retained minimal preamble for the Stacks Project chapter sites-cohomology.tex (Cohomology on Sites).
% Verbatim excerpt of the chapter preamble: Xy-pic setup (native preamble lines 14-18) and the native
% theorem declarations with the mathematical macros (native preamble lines 154-217). Copied mechanically;
% nothing added, removed or redefined. The wrapper class and the amsmath/amssymb/amsthm/graphicx base
% packages are supplied by the assembler, so the class/externaldocument part of the native preamble is omitted.
% For commutative diagrams we use Xy-pic
\usepackage[all]{xy}

% We use 2cell for 2-commutative diagrams.
\xyoption{2cell}
% Theorem environments.
%
\theoremstyle{plain}
\newtheorem{theorem}[subsection]{Theorem}
\newtheorem{proposition}[subsection]{Proposition}
\newtheorem{lemma}[subsection]{Lemma}

\theoremstyle{definition}
\newtheorem{definition}[subsection]{Definition}
\newtheorem{example}[subsection]{Example}
\newtheorem{exercise}[subsection]{Exercise}
\newtheorem{situation}[subsection]{Situation}

\theoremstyle{remark}
\newtheorem{remark}[subsection]{Remark}
\newtheorem{remarks}[subsection]{Remarks}

\numberwithin{equation}{subsection}

% Macros
%
\def\lim{\mathop{\mathrm{lim}}\nolimits}
\def\colim{\mathop{\mathrm{colim}}\nolimits}
\def\Spec{\mathop{\mathrm{Spec}}}
\def\Hom{\mathop{\mathrm{Hom}}\nolimits}
\def\Ext{\mathop{\mathrm{Ext}}\nolimits}
\def\SheafHom{\mathop{\mathcal{H}\!\mathit{om}}\nolimits}
\def\SheafExt{\mathop{\mathcal{E}\!\mathit{xt}}\nolimits}
\def\Sch{\mathit{Sch}}
\def\Mor{\mathop{\mathrm{Mor}}\nolimits}
\def\Ob{\mathop{\mathrm{Ob}}\nolimits}
\def\Sh{\mathop{\mathit{Sh}}\nolimits}
\def\NL{\mathop{N\!L}\nolimits}
\def\CH{\mathop{\mathrm{CH}}\nolimits}
\def\proetale{{pro\text{-}\acute{e}tale}}
\def\etale{{\acute{e}tale}}
\def\QCoh{\mathit{QCoh}}
\def\Ker{\mathop{\mathrm{Ker}}}
\def\Im{\mathop{\mathrm{Im}}}
\def\Coker{\mathop{\mathrm{Coker}}}
\def\Coim{\mathop{\mathrm{Coim}}}

% Boxtimes
%
\DeclareMathSymbol{\boxtimes}{\mathbin}{AMSa}{"02}

%
% Macros for moduli stacks/spaces
%
\def\QCohstack{\mathcal{QC}\!\mathit{oh}}
\def\Cohstack{\mathcal{C}\!\mathit{oh}}
\def\Spacesstack{\mathcal{S}\!\mathit{paces}}
\def\Quotfunctor{\mathrm{Quot}}
\def\Hilbfunctor{\mathrm{Hilb}}
\def\Curvesstack{\mathcal{C}\!\mathit{urves}}
\def\Polarizedstack{\mathcal{P}\!\mathit{olarized}}
\def\Complexesstack{\mathcal{C}\!\mathit{omplexes}}
% \Pic is the operator that assigns to X its picard group, usage \Pic(X)
% \Picardstack_{X/B} denotes the Picard stack of X over B
% \Picardfunctor_{X/B} denotes the Picard functor of X over B
\def\Pic{\mathop{\mathrm{Pic}}\nolimits}
\def\Picardstack{\mathcal{P}\!\mathit{ic}}
\def\Picardfunctor{\mathrm{Pic}}
\def\Deformationcategory{\mathcal{D}\!\mathit{ef}}
\usepackage{cleveref}
\crefname{originalremark}{original remark}{original remarks}
\Crefname{originalremark}{Original remark}{Original remarks}
\makeatletter
\newlabel{remark-cup-product}{{0B6C}{source}{}{sourcepointer-1}{}}
\newlabel{remark-cup-product@cref}{{[originalremark][0][]0B6C}{[1][0][]source}{}{}{}}
\makeatother
\title{Associativity of the relative cup product for a morphism of ringed topoi}
\author{The Stacks Project Authors}
\date{stacks-project:sites-cohomology.tex:local-snapshot-20261011; GNU FDL 1.2 or later}
\begin{document}
\maketitle
\noindent\textit{Editorial note.} All delivered mathematical text is the author's own, copied from the pinned Stacks Project chapter file sites-cohomology.tex; no statement or proof step was rewritten, reordered or removed.\par\medskip
\noindent\textit{Editorial note.} The proof is the author's complete argument for the commutativity of the diagram: the transposition is displayed and the two sides are identified in the same display.\par\medskip
\noindent\textit{Editorial note.} The single labelled reference, which occurs inside the counted statement, is carried in assembly.external\_references as a non-load-bearing navigation pointer to the same chapter.\par\medskip
\noindent\textit{Editorial note.} Cross-chapter restatement: the same label occurs in chapter cohomology for a ringed space, and where applicable in chapter more-algebra for a ring. Those statements are the special cases of this ringed-site statement, so per the lane-s4 rule only this more general occurrence is counted here; the special-case twins are to be left uncounted by the lanes holding those chapters.\par\medskip
\noindent\textit{Editorial note.} The preamble is the chapter's own retained minimal preamble (sources/stacks-sites-cohomology/retained-preamble.tex), a verbatim excerpt of the native preamble: the Xy-pic setup and the native theorem declarations with the mathematical macros. The Xy-pic diagram inside the counted statement is delivered as the author wrote it.\par\medskip
\noindent\textit{Allowed standard prerequisites (editorial).}\begin{itemize}
\item Relative cup product for a morphism of ringed topoi. Same chapter, the morphism Rf\_*K tensor\textasciicircum{}L Rf\_*L -> Rf\_*(K tensor\textasciicircum{}L L) named in the counted statement. Supplies the two arrows whose composite is asserted to be independent of the order of bracketing.
\item Derived pullback Lf\textasciicircum{}* and the projection-style adjunction with Rf\_*. Same chapter, the adjunction between Lf\textasciicircum{}* and Rf\_* on derived categories of modules. Supplies the passage to adjoints with which the proof replaces the comparison of two composites by the comparison of two maps in D(O).
\item Monoidality of the derived pullback. Same chapter, the identification Lf\textasciicircum{}*(A tensor\textasciicircum{}L B) = Lf\textasciicircum{}*A tensor\textasciicircum{}L Lf\textasciicircum{}*B used in the displayed computation. Supplies the first equality in the displayed chain, which produces the common adjoint map.
\end{itemize}\medskip
\ifdefined\hypertarget\hypertarget{sourcepointer-1}{}\fi
\noindent\textit{Source reference.} Original remark 0B6C: sources/stacks-sites-cohomology/pinned-original-source.tex, source lines 4193--4193. Same-chapter navigation pointer occurring inside the counted statement, where the author names the relative cup product of that remark. The pointer only resolves the reference for the reader; the remark is not reproduced inside the counted span, and no bibliography citation is involved.\par\medskip
\begin{lemma}
\label{lemma-cup-product-associative}
Let $f : (\Sh(\mathcal{C}), \mathcal{O}) \to (\Sh(\mathcal{C}'), \mathcal{O}')$
be a morphism of ringed topoi. The relative cup product of
Remark \ref{remark-cup-product} is associative in the sense that
the diagram
$$
\xymatrix{
Rf_*K \otimes_{\mathcal{O}'}^\mathbf{L}
Rf_*L \otimes_{\mathcal{O}'}^\mathbf{L}
Rf_*M \ar[r] \ar[d] &
Rf_*(K \otimes_\mathcal{O}^\mathbf{L} L)
\otimes_{\mathcal{O}'}^\mathbf{L} Rf_*M \ar[d] \\
Rf_*K \otimes_{\mathcal{O}'}^\mathbf{L}
Rf_*(L \otimes_\mathcal{O}^\mathbf{L} M) \ar[r] &
Rf_*(K \otimes_\mathcal{O}^\mathbf{L} 
L \otimes_\mathcal{O}^\mathbf{L} M)
}
$$
is commutative in $D(\mathcal{O}')$ for all $K, L, M$ in $D(\mathcal{O})$.
\end{lemma}

\begin{proof}
Going around either side we obtain the map adjoint to the obvious map
\begin{align*}
Lf^*(Rf_*K \otimes_{\mathcal{O}'}^\mathbf{L}
Rf_*L \otimes_{\mathcal{O}'}^\mathbf{L}
Rf_*M) & =
Lf^*(Rf_*K) \otimes_\mathcal{O}^\mathbf{L}
Lf^*(Rf_*L) \otimes_\mathcal{O}^\mathbf{L}
Lf^*(Rf_*M) \\
& \to
K \otimes_\mathcal{O}^\mathbf{L} 
L \otimes_\mathcal{O}^\mathbf{L} M
\end{align*}
in $D(\mathcal{O})$.
\end{proof}

\end{document}
